\documentclass[10pt,a4paper]{amsart}
\usepackage[margin=19mm]{geometry}
\usepackage{amsmath,amssymb,mathtools}
\usepackage[scr=rsfs,cal=euler]{mathalfa}
\usepackage{xfrac}
\usepackage[unicode,hidelinks,bookmarks=false]{hyperref}
\newcommand{\ie}{i.e.\ }
\newcommand{\cf}{cf.\ }
\newcommand{\resp}{respectively\ }
\newcommand{\Subsection}{\subsection}
\providecommand{\nobreakdash}{\nobreak\hbox{-}}
\setlength{\emergencystretch}{2em}
\multlinegap0pt
\usepackage{amssymb}
\usepackage{float}
\newtheorem{proposition}{Proposition}
\title{On a sequence of Kimberling and its relationship to the Tribonacci word}
\author{Lubom\'ira Dvo\v{r}\'akov\'a, Edita Pelantov\'a,, Jeffrey Shallit}
\date{Source: arXiv:2510.11318v5 (CC BY 4.0)}
\begin{document}
\maketitle

\noindent\emph{Source and licence.} On a sequence of Kimberling and its relationship to the Tribonacci word. Version arXiv:2510.11318v5 (CC BY 4.0), distributed by arXiv under the Creative Commons Attribution 4.0 International licence, http://creativecommons.org/licenses/by/4.0/, as recorded in the arXiv licence metadata for this exact version and shown on the abstract page https://arxiv.org/abs/2510.11318v5. Full licence notice: \texttt{licenses/arxiv-2510.11318-CC-BY-4.0.txt}.\par\smallskip

\noindent\textit{Editorial scope.} Counted unit: the article's proposition proposition (source0.tex lines 187--190). The article's complete original proof of it is delivered with it (source0.tex lines 192--260, one \texttt{proof} environment of 69 lines). No mathematics is written, repaired or re-derived: the statement and the proof are the article's own lines, in the article's own notation, and no retained line is substituted or re-flowed.  No further source-local context is needed: every cross-reference of every retained line resolves inside the two delivered spans.  Nothing was omitted from any retained span, so the package carries no omission record.  No retained line contains a citation, so the wrapper carries no bibliography and no bibliography entry.  No retained line contains a figure environment, an image-inclusion command or drawing code, so no image is delivered and the package declares no asset dependency.  Editorial formatting only, in addition: an AMS \texttt{amsart} preamble that restates the article's own declarations and macros used by the retained lines, namely the environments \texttt{proposition} on separate counters, together with the standard packages those lines need.  The retained environments print in this document as Proposition 1 in order of appearance, whereas the article prints the same items with the numbering of its own full document.  The complete native source of the article as distributed in the arXiv e-print for this exact version is bundled unchanged as \texttt{sources/186953/source0.tex}.
\begin{proposition}
Let $c_0 = 2$, $c_1 = 4$, $c_2 = 6$, $c_3 = 10$, and
define $c_i = 2c_{i-1} - c_{i-4}$ for $i \geq 4$.  Then $|B_i|= c_i$ for $i \geq 0$.
\end{proposition}

\begin{proof}
A simple induction now shows that, for $n \geq 1$, the word $B_i$ starts with $01$
and contains no occurrences of either $000$ or $111$.  Therefore $B_i$ can be
factorized uniquely as a concatenation of the single initial $0$ and blocks of the
form $1$, $10$, and $100$.  For $x \in \{ 1, 10, 100 \}$, define $N_x (i)$ to be
the number of occurrences of the block $x$ in this
factorization of $B_i$.  Then clearly
$|B_i| = 1 + N_1(i) + 2N_{10} (i) + 3N_{100} (i)$.

Define a mapping $\xi$ on the blocks $x$ that obeys Kimberling's inflation 
rules, and hence sends $1$ to $10$, $10$ to $100$, and
$100$ to $100101$.  Write $B_i = 0 x_1 x_2 \cdots x_k$, where each $x_j \in \{ 1, 10, 100 \}$, $1 \leq j \leq k$. By checking what happens at the boundaries, we see that $B_{i+1} = 0 \xi(x_1) \cdots \xi(x_k)$.  

We now claim that the following relations hold for $i \geq 2$:
\begin{align}
N_1 (i) &= N_{100} (i-1) \label{eqn1} \\
N_{10} (i) &= N_{100} (i-1) + N_1 (i-1) \label{eqn10} \\
N_{100} (i) &= N_{10} (i-1) + N_{100} (i-1). \label{eqn100}
\end{align}
These follow immediately by looking at the image of each block above.

We now show that these identities are enough to directly prove the following (no induction needed!):
\begin{align}
N_1 (i) &= N_1(i-1) + N_1(i-2) + N_1(i-3) \label{en1}\\
N_{10} (i) &= N_{10}(i-1) + N_{10}(i-2) + N_{10}(i-3) \label{n10} \\
N_{100} (i) &= N_{100}(i-1) + N_{100}(i-2) + N_{100} (i-3) \label{n100} 
\end{align}
for $i \geq 4$. 

Proof of Eq.~\eqref{en1}:
\begin{align*}
N_1 (i) &= N_{100} (i-1) \quad \text{(by Eq.~\eqref{eqn1})}\\
&= N_{10} (i-2) + N_{100}(i-2)  \quad \text{(by Eq.~\eqref{eqn100})}\\
&= N_{10} (i-2) + N_1(i-1)  \quad \text{(by Eq.~\eqref{eqn1})} \\
&= N_{100} (i-3) + N_1(i-3) + N_1(i-1) \quad \text{(by Eq.~\eqref{eqn100})} \\
&= N_1(i-2) + N_1(i-3) + N_1(i-1) \quad \text{(by Eq.~\eqref{eqn1})} .
\end{align*}

Proof of Eq.~\eqref{n10}:
\begin{align*}
N_{100} (i) &= N_{10} (i-1) + N_{100} (i-1) \quad \text{(by Eq.~\eqref{eqn100})} \\
&= N_{100} (i-2) + N_1 (i-2) + N_{100} (i-1) \quad \text{(by Eq.~\eqref{eqn10})} \\
&= N_{100} (i-2) + N_{100} (i-3) + N_{100} (i-1) \quad \text{(by Eq.~\eqref{eqn1})}.
\end{align*}

Proof of Eq.~\eqref{n100}:
\begin{align}
N_{10} (i) &= N_{100} (i-1) + N_1 (i-1) \quad \text{(by Eq.~\eqref{eqn10})}  \nonumber\\
&= N_{10} (i-2) + N_{100} (i-2) + N_1(i-1) \quad \text{(by Eq.~\eqref{eqn100})} \nonumber \\
&= N_{10} (i-2) + (N_{10} (i-1) - N_1 (i-2)) + N_1 (i-1) \quad \text{(by Eq.~\eqref{eqn10})} \label{done} .
\end{align}

We also have
\begin{align*}
N_1(i-1) - N_1 (i-2) &= N_{100} (i-2) - N_{100} (i-3) \quad \text{(by Eq.~\eqref{eqn1})} \\
&= N_{10} (i-3) \quad \text{(by Eq.~\eqref{eqn100})},
\end{align*}
and substituting into Eq.~\eqref{done} gives the desired result for $N_{10} (i)$.

Now from above we know that $|B_i| - 1 = N_1 (i) + 2 N_{10} (i) + 3 N_{100} (i)$.  Since each of the
sequences $(N_1 (i))_i$, $(N_{10} (i))_i$, and $(N_{100} (i))_i$ 
on the right-hand side is annihilated by the shift polynomial $X^3 - X^2 - X - 1$, so is their linear
combination  $(N_1 (i) + 2 N_{10} (i) + 3 N_{100} (i))_i$.  And since $X-1$ annihilates the
constant sequence $1$, the sequence $(|B_i|)_i$ is annihilated by the product of
$X-1$ and $X^3 - X^2 - X - 1$, which is $X^4 - 2X^3 + 1$.   In other words,
$|B_{i}| = 2|B_{i-3}| - |B_{i-4}|$ for $i \geq 4$.  After comparing the initial conditions
$i = 0,1,2,3$, it now follows that
$|B_i| = c_i$ for all $i \geq 0$.
\end{proof}

\end{document}
